% 图11.1：原始四面板矢量数据、置信区间及截断箭头均原样保留；只替换英文标签。
\begin{tikzpicture}[x=.45mm,y=-.45mm]
 \node[anchor=north west,inner sep=0pt] at (0,0)
  {\includegraphics[width=150.3mm,height=96.3mm]{figures/ch11_sieve_data.pdf}};
 \node[font=\small,rotate=90] at (5.9503,45.9548) {累积发生率};
 \node[anchor=west,font=\fontsize{7.2}{7.7}\selectfont] at (35.9112,17.3084) {安慰剂};
 \node[anchor=west,font=\fontsize{7.2}{7.7}\selectfont] at (35.9112,24.2791) {疫苗};
 \node[font=\small] at (87.8391,6.5254) {氨基酸 169 位点匹配};
 \node[font=\small,rotate=90] at (173.3625,45.9543) {累积发生率};
 \node[anchor=west,font=\fontsize{7.2}{7.7}\selectfont] at (203.9068,17.3085) {安慰剂};
 \node[anchor=west,font=\fontsize{7.2}{7.7}\selectfont] at (203.9068,24.2791) {疫苗};
 \node[font=\small] at (255.8328,6.5254) {氨基酸 169 位点不匹配};
 \node[font=\small] at (87.8429,208.2375) {末次疫苗接种后的月数};
 \node[font=\small,rotate=90] at (5.9508,147.5245) {疫苗效力};
 \node[anchor=west,font=\fontsize{7.2}{7.7}\selectfont] at (40.5308,114.8452) {169 位点匹配};
 \node[anchor=west,font=\fontsize{7.2}{7.7}\selectfont] at (40.5308,121.8159) {169 位点不匹配};
 \node[font=\small] at (255.8344,208.2375) {末次疫苗接种后的月数};
 \node[font=\small,rotate=90] at (173.3615,147.5202) {筛选效应};
\end{tikzpicture}
